\section*{Chapitre \ref{ch:rs:overfitting} --- Surapprentissage}

\begin{solution}{exo:rs:overfitting:1}
$\omega = 380/1\,001 = 0.380$, $\lambda = \ln(0.380/0.620) = -0.49$~: sous la médiane, donc la partition compte dans la PBO.
\end{solution}

\begin{solution}{exo:rs:overfitting:2}
$\binom{10}{2}\times 2/10 = 45 \times 2/10 = 9$ chemins~; $\binom{6}{3} \times 3/6 = 10$.
\end{solution}

\begin{solution}{exo:rs:overfitting:3}
$2\ln 100/1^2 = 9.2$ ans.
\end{solution}

\begin{solution}{exo:rs:overfitting:4}
Une étiquette observée le jour $d$ couvre les jours $d$ à $d + 9$. Les étiquettes de test couvrent les jours 200 à 308. La \omterm{def:rs:overfitting:purging}{purge} retire les jours d'entraînement dont les étiquettes touchent cet intervalle~:
191 à 199 (leurs étiquettes atteignent le jour 200 et les suivants) et 300 à 308 (leurs étiquettes commencent à l'intérieur de cet intervalle). Un \omterm{def:rs:overfitting:purging}{embargo} de 5 jours retire aussi les jours 309 à 313.
\end{solution}

\begin{solution}{exo:rs:overfitting:5}
Les systèmes corrélés partagent l'essentiel de leur bruit~; le gagnant \omterm{def:rs:overfitting:overfitting}{dans l'échantillon} est celui dont la part idiosyncratique s'ajuste le mieux à la moitié de l'échantillon. Les moitiés
étant tirées symétriquement, la même part idiosyncratique tend à moins bien ajuster l'autre moitié que la part commune ne le fait pour les autres, si bien que le gagnant tombe sous
la médiane plus souvent qu'à son tour~: la sélection est anticorrélée au rang \omterm{def:rs:overfitting:overfitting}{hors échantillon} (les pentes de $-0.89$ et $-0.91$). Avec des stratégies indépendantes et aucune
compétence, la PBO avoisine un demi.
\end{solution}

\begin{solution}{exo:rs:overfitting:6}
La tendance est commune à tous les systèmes~: elle les soulève tous à peu près de la même quantité et ne permet pas de les distinguer. Ce qui les distingue \omterm{def:rs:overfitting:overfitting}{dans l'échantillon} est
du bruit, qui ne persiste pas~; le gagnant d'une recherche sur du bruit superposé à une tendance commune obtient la tendance moins sa pénalité de sélection (0.14 contre une
moyenne de 0.26).
\end{solution}

\begin{solution}{exo:rs:overfitting:7}
\texttt{rs\_overfit.pbo\_hundred}~: 0.69 pour les 100 premiers systèmes à tendance, 0.52 pour 100 stratégies indépendantes de pur bruit. Des stratégies indépendantes sans compétence donnent une
PBO proche d'un demi, comme il se doit pour une sélection parmi du pur bruit~; des stratégies corrélées la font dépasser ce niveau, car l'edge du gagnant est la part de son bruit que les autres ne partagent pas.
\end{solution}

\begin{solution}{exo:rs:overfitting:8}
Des plis mélangés placent des étiquettes qui se chevauchent et des variables explicatives sérialement corrélées de part et d'autre de chaque découpage~: le modèle est évalué sur des jours dont il a, en
effet, déjà vu les étiquettes. Sur le bruit du chapitre, les plis mélangés donnent un $R^2$ de 0.90, et tout schéma respectant l'ordre temporel donne une valeur négative. Utiliser des plis contigus,
purger les étiquettes qui se chevauchent, appliquer un \omterm{def:rs:overfitting:purging}{embargo} après les plis de test, et calibrer sur des plis qui précèdent la période de test.
\end{solution}

\begin{solution}{pb:rs:overfitting:1}
\begin{enumerate}
\item Bruit~: 0.44 \omterm{def:rs:overfitting:overfitting}{dans l'échantillon}, $-0.32$ \omterm{def:rs:overfitting:overfitting}{hors échantillon}. Tendance~: 0.77 et 0.14.
\item Bruit~: médiane \omterm{def:rs:overfitting:overfitting}{dans l'échantillon} 0.17, moyenne \omterm{def:rs:overfitting:overfitting}{hors échantillon} 0.20. Tendance~: 0.50 et 0.26.
\item 0.77~; les systèmes sont fortement corrélés (leurs ratios de Sharpe s'étalent sur 0.076), si bien que la recherche revient à peu d'essais indépendants, et le meilleur attendu
compte tenu de cette dispersion, se situe 0.25 au-dessus du niveau commun.
\item La tendance soulève chaque système (médiane 0.50). La recherche ajoute un classement qui transporte une part d'information sur l'ensemble des systèmes (corrélation 0.33,
les cinquante meilleurs obtiennent 0.29), mais elle choisit son gagnant unique surtout pour du bruit.
\item À partir des sommes par blocs et des sommes de carrés~: chaque partition est un ensemble de blocs, si bien que ses moyennes et ses variances sont des produits matriciels.
\item Bruit~: PBO 0.82, pente $-0.89$. Tendance~: 0.74, pente $-0.91$.
\item La sélection fait pire que le choix au hasard~: le gagnant \omterm{def:rs:overfitting:overfitting}{dans l'échantillon} se situe plus souvent qu'à son tour sous la médiane \omterm{def:rs:overfitting:overfitting}{hors échantillon}.
\item Proche de zéro~: il gagne \omterm{def:rs:overfitting:overfitting}{dans l'échantillon} et reste en tête \omterm{def:rs:overfitting:overfitting}{hors échantillon} dans presque toutes les partitions (le test de la construction trouve moins de 1~\%).
\item 42.4, 10.6 et 3.3 ans.
\item L'erreur type d'un ratio de Sharpe décroît comme la racine carrée de l'historique~; pour diviser par deux l'erreur visée, il faut
quatre fois plus d'historique.
\item Mélangés 0.90, contigus $-0.63$, purgés $-0.57$, purgés et avec \omterm{def:rs:overfitting:purging}{embargo} $-0.93$.
\item Les jours voisins, dont les étiquettes partagent des périodes, tombent dans les deux ensembles~; la \omterm{def:rs:overfitting:purging}{purge} retire les jours d'entraînement dont les étiquettes chevauchent celles du test, l'\omterm{def:rs:overfitting:purging}{embargo}
retire les jours juste après l'ensemble de test, dont les variables explicatives portent encore son information.
\item $-0.04$ sur du bruit, 0.29 avec la tendance.
\item \textbf{Résultat nommé.} La PBO de la recherche vaut 0.82 sur du bruit et 0.74 avec une vraie tendance~; le ratio de Sharpe \omterm{def:rs:overfitting:overfitting}{hors échantillon} du système retenu vaut $-0.32$ et 0.14, en dessous
de la moyenne des systèmes (0.20 et 0.26 dans chacun des deux mondes), tandis que la resélection annuelle en walk-forward obtient $-0.04$ et 0.29. L'espérance déflatée du gagnant est la
moyenne des systèmes, et non son propre chiffre \omterm{def:rs:overfitting:overfitting}{dans l'échantillon}, 0.44 ou 0.77.
\item L'\omterm{def:rs:overfitting:walkforward}{analyse glissante} (walk-forward), ou la moyenne du plateau, qui a rapporté la tendance sans payer le choix d'un point.
\item Les mille couples de paramètres comptés comme essais, la période de données, la règle de sélection, la PBO, la \omterm{def:rs:overfitting:minbtl}{longueur minimale du backtest} et le ratio de Sharpe déflaté, avant tout
examen \omterm{def:rs:overfitting:overfitting}{hors échantillon}.
\item Une raison de la tendance (pression de couverture, diffusion lente) prédit que la classe large des systèmes doit fonctionner, et non une fenêtre particulière, et désigne le
plateau plutôt que le pic.
\item Une décote du ratio de Sharpe publié pour le nombre de tests, un seul chiffre à placer à côté de la probabilité PBO.
\item De nombreux chemins \omterm{def:rs:overfitting:overfitting}{hors échantillon} au lieu d'un seul, de sorte que la performance de la procédure soit une distribution.
\item Choisir, parmi de nombreux \omterm{def:rs:vectorised-backtests:backtest}{backtests} sur un historique limité, celui dont le succès tient à l'échantillon plutôt qu'à la stratégie.
\end{enumerate}
\end{solution}

\begin{solution}{iq:rs:overfitting:1}
Consigner les 500 essais~; calculer le ratio de Sharpe déflaté du meilleur (livre~4, chapitre~12) et la PBO de la recherche par CSCV~; vérifier la \omterm{def:rs:overfitting:minbtl}{longueur minimale du backtest} pour
500 essais à un ratio de Sharpe de 2~; examiner le voisinage du meilleur (plateau ou pic)~; trouver un mécanisme~; puis tester sur des données non touchées par la recherche, ou trader
à petite échelle (chapitre~21).
\end{solution}

\begin{solution}{iq:rs:overfitting:2}
Les observations sont sérialement corrélées et les étiquettes se chevauchent, si bien que des plis mélangés font fuiter l'information du test vers l'entraînement. Utiliser des plis contigus avec \omterm{def:rs:overfitting:purging}{purge} des
étiquettes qui se chevauchent, un \omterm{def:rs:overfitting:purging}{embargo} après chaque pli de test, une \omterm{def:rs:overfitting:walkforward}{analyse glissante} (walk-forward) pour la procédure, et une \omterm{def:rs:overfitting:purging}{validation croisée combinatoire purgée} (CPCV) pour une distribution de chemins.
\end{solution}

\begin{solution}{iq:rs:overfitting:3}
La PBO est la probabilité que le gagnant d'une recherche \omterm{def:rs:overfitting:overfitting}{dans l'échantillon} se situe sous la médiane \omterm{def:rs:overfitting:overfitting}{hors échantillon}. La CSCV découpe l'historique en blocs, prend chaque
partition symétrique en moitiés \omterm{def:rs:overfitting:overfitting}{dans l'échantillon} et \omterm{def:rs:overfitting:overfitting}{hors échantillon}, choisit le gagnant \omterm{def:rs:overfitting:overfitting}{dans l'échantillon} et le classe \omterm{def:rs:overfitting:overfitting}{hors échantillon}~; la part des partitions dont le rang est sous la médiane
estime la PBO.
\end{solution}

\begin{solution}{iq:rs:overfitting:4}
Environ $2\ln 1\,000 = 13.8$ ans par l'approximation, 10.6 avec le maximum attendu exact d'essais indépendants~; des essais corrélés en demandent moins, et une réponse
digne de confiance en demande davantage (pour que le niveau de hasard reste nettement sous la cible, et non égal à celle-ci).
\end{solution}

\begin{solution}{iq:rs:overfitting:5}
Un \omterm{def:rs:the-research-process:log}{journal de recherche} où chaque essai est compté~; des protocoles fixés avant les résultats~; un échantillon de réserve gardé par un tiers~; un examen des recherches par leur PBO et leurs statistiques déflatées~;
une valorisation des mécanismes et des résultats négatifs~; des limites au nombre de variantes qu'un projet peut tester par année de données.
\end{solution}

\begin{solution}{iq:rs:overfitting:6}
Les estimations du ratio de Sharpe nul sur $y$ ans sont approximativement normales, de variance $1/y$~; le maximum attendu de $N$ d'entre elles vaut environ $\sqrt{2\ln N}/\sqrt y$~; l'égaler à la
cible $s$ donne $y = 2\ln N/s^2$~: logarithmique en fonction du \omterm{def:rs:the-research-process:log}{nombre d'essais}, en raison inverse du carré de la cible.
\end{solution}
